\documentclass[10pt, letterpaper]{article}

% Packages:
\usepackage[
    ignoreheadfoot,
    top=0.4 cm,
    bottom=0.4 cm,
    left=0.8 cm,
    right=0.8 cm,
    footskip=0.2 cm,
]{geometry}
\usepackage{titlesec}
\usepackage{tabularx}
\usepackage{array}
\usepackage[dvipsnames]{xcolor}
\definecolor{primaryColor}{RGB}{0, 0, 0}
\usepackage{enumitem}
\usepackage{amsmath}
\usepackage[
    pdftitle={Tanmay Bhuskute's CV},
    pdfauthor={Tanmay Bhuskute},
    pdfcreator={LaTeX with RenderCV},
    colorlinks=true,
    urlcolor=primaryColor
]{hyperref}
\usepackage[pscoord]{eso-pic}
\usepackage{calc}
\usepackage{bookmark}
\usepackage{lastpage}
\usepackage{changepage}
\usepackage{paracol}
\usepackage{ifthen}
\usepackage{needspace}
\usepackage{iftex}

\ifPDFTeX
    \input{glyphtounicode}
    \pdfgentounicode=1
    \usepackage[T1]{fontenc}
    \usepackage[utf8]{inputenc}
    \usepackage{lmodern}
\fi

\usepackage{XCharter}

% Some settings:
\raggedright
\AtBeginEnvironment{adjustwidth}{\partopsep0pt}
\pagestyle{empty}
\setcounter{secnumdepth}{0}
\setlength{\parindent}{0pt}
\setlength{\topskip}{0pt}
\setlength{\columnsep}{0.15cm}
\pagenumbering{gobble}

% ── Section title ──
\titleformat{\section}{\needspace{2\baselineskip}\bfseries\large}{}{0pt}{}[\vspace{1pt}\titlerule]
\titlespacing{\section}{-1pt}{0.15 cm}{0.08 cm}

\renewcommand\labelitemi{$\vcenter{\hbox{\small$\bullet$}}$}

% ── Bullet lists ──
\newenvironment{highlights}{
    \begin{itemize}[
        topsep=0.10 cm,
        parsep=0.06 cm,
        partopsep=0pt,
        itemsep=0.06 cm,
        leftmargin=12pt
    ]
}{
    \end{itemize}
}

\newenvironment{highlightsforbulletentries}{
    \begin{itemize}[
        topsep=0.10 cm,
        parsep=0.06 cm,
        partopsep=0pt,
        itemsep=0.04 cm,
        leftmargin=12pt
    ]
}{
    \end{itemize}
}

\newenvironment{onecolentry}{
    \begin{adjustwidth}{0 cm + 0.00001 cm}{0 cm + 0.00001 cm}
}{
    \end{adjustwidth}
}

\newenvironment{twocolentry}[2][]{
    \onecolentry
    \def\secondColumn{#2}
    \setcolumnwidth{\fill, 4.5 cm}
    \begin{paracol}{2}
}{
    \switchcolumn \raggedleft \secondColumn
    \end{paracol}
    \endonecolentry
}

\newenvironment{threecolentry}[3][]{
    \onecolentry
    \def\thirdColumn{#3}
    \setcolumnwidth{, \fill, 4.5 cm}
    \begin{paracol}{3}
    {\raggedright #2} \switchcolumn
}{
    \switchcolumn \raggedleft \thirdColumn
    \end{paracol}
    \endonecolentry
}

\newenvironment{header}{
    \setlength{\topsep}{0pt}\par\kern\topsep\centering\linespread{0.6}
}{
    \par\kern\topsep
}

\let\hrefWithoutArrow\href

\begin{document}
    \newcommand{\AND}{\unskip
        \cleaders\copy\ANDbox\hskip\wd\ANDbox
        \ignorespaces
    }
    \newsavebox\ANDbox
    \sbox\ANDbox{$|$}

    \begin{header}
        \fontsize{16 pt}{24 pt}\selectfont Tanmay Bhuskute

        \vspace{3 pt}

        \normalsize
        \mbox{ +1 (732) 858-3027 }%
        \kern 1.0 pt%
        \AND%
        \kern 1.0 pt%
        \mbox{ tanmay.v.bhuskute@gmail.com }%
        \kern 1.0 pt%
        \AND%
        \kern 1.0 pt%
        \mbox{ github.com/Tanmayb05 }%
        \kern 1.0 pt%
        \AND%
        \kern 1.0 pt%
        \mbox{ linkedin.com/in/tanmay-bhuskute }%
        \kern 1.0 pt%
        \AND%
        \kern 1.0 pt%
        \mbox{ Livermore, CA }%
    \end{header}

    \vspace{2 pt}

%═══════════════════════════════════════
\section{SUMMARY}
%═══════════════════════════════════════
\begin{onecolentry}
Software Engineer and M.S. Computer Science candidate (Dec 2026) with production experience writing Python software for data acquisition, processing, and analysis. Skilled in building experimental harnesses, reproducible data pipelines, and full-stack analysis tools under rigorous coding standards, and in collaborating with scientists and technical staff on research goals.
\end{onecolentry}

%═══════════════════════════════════════
\section{EDUCATION}
%═══════════════════════════════════════
\begin{twocolentry}{
    Aug 2024 - Dec 2026
}
    \textbf{Master of Science in Computer Science} \\
    University of Cincinnati
\end{twocolentry}
\begin{onecolentry}
Relevant Coursework: Advanced Algorithms, Deep Learning, Intelligent Data Analysis.
\end{onecolentry}
\vspace{0.15 cm}

\begin{twocolentry}{
    Aug 2018 - Jul 2022
}
    \textbf{Bachelor of Engineering in Computer Engineering} \\
    AISSMS College of Engineering
\end{twocolentry}
\begin{onecolentry}
Relevant Coursework: Data Mining, Data Analytics, Object-Oriented Programming.
\end{onecolentry}
%═══════════════════════════════════════
\section{TECHNICAL SKILLS}
%═══════════════════════════════════════
\begin{onecolentry}
\textbf{Programming:} Python, Java, C/C++, JavaScript, TypeScript, SQL, Go, Bash.
\end{onecolentry}
\vspace{0.06 cm}
\begin{onecolentry}
\textbf{Data Acquisition \& Analysis:} pandas, NumPy, scikit-image, PySpark, Delta Lake, Matplotlib, OpenCV, SQL analytics.
\end{onecolentry}
\vspace{0.06 cm}
\begin{onecolentry}
\textbf{Web \& Backend:} FastAPI, Flask, Node.js, React, REST APIs, PostgreSQL, MySQL, Redis.
\end{onecolentry}
\vspace{0.06 cm}
\begin{onecolentry}
\textbf{Developer Tools:} Git, GitLab CI/CD, Jenkins, Linux, macOS, Docker, Kubernetes, AWS, SonarQube, Pytest, JUnit.
\end{onecolentry}
\vspace{0.06 cm}
\begin{onecolentry}
\textbf{AI \& ML:} Gemini, Vertex AI, PyTorch, TensorFlow, scikit-learn, LangChain, AI coding tools (Claude, Copilot).
\end{onecolentry}

%═══════════════════════════════════════
\section{WORK EXPERIENCE}
%═══════════════════════════════════════
\begin{twocolentry}{Feb 2026 -- Present}
\textbf{Graduate Researcher} \,|\, University of Cincinnati
\end{twocolentry}
\begin{highlights}
  \item Built a config-driven experimental harness in \textbf{Python} for a two-stage LLM pipeline reproducing Android GUI bugs from screen recordings, with staged execution (\texttt{--stage 1/2/all}) and dry-run validation for reproducible evaluation runs.
  \item Designed a flat, versioned run-artifact schema keeping every experimental trial reproducible and auditable, akin to a detailed lab notebook; evaluated against \textbf{ViBR} baseline across \textbf{53 runs} on \textbf{20+ apps}, showing \textbf{25.7x latency reduction} (39.4s vs 1012.1s avg).
  \item Prepared progress reports and comparative results for faculty mentors, isolating method effects from model choice via a provider-agnostic LLM abstraction.
\end{highlights}

\vspace{0.15 cm}

\begin{twocolentry}{Jul 2022 -- Jun 2024}
\textbf{Software Development Engineer} \,|\, Siemens Industry Software
\end{twocolentry}
\begin{highlights}
  \item Wrote \textbf{Python} (Boto3) automation for infrastructure data collection and deployment workflows, replacing legacy Ruby scripts and reducing deployment time by \textbf{40\%} across AWS environments.
  \item Executed \textbf{44 production environment upgrades} in 4 months, diagnosing and documenting live issues across Linux, EC2, databases, and services in a fast-paced operational setting.
  \item Served as \textbf{Quality Coach and SonarQube Specialist}, enforcing rigorous coding standards and CI/CD quality gates that raised automated code coverage from \textbf{35\% to 78\%}.
  \item Mentored \textbf{5 software engineers} and interfaced across engineering teams through structured onboarding, technical documentation, and standardized development practices.
\end{highlights}

\vspace{0.15 cm}

\begin{twocolentry}{Sep 2021 -- Dec 2021}
\textbf{Software Developer Intern} \,|\, Government of Maharashtra, Water Resource Department
\end{twocolentry}
\begin{highlights}
  \item Built and deployed a full-stack \textbf{React and MySQL} application that digitized departmental record management, automated manual data workflows, and improved operational efficiency by \textbf{40\%}.
\end{highlights}

%═══════════════════════════════════════
\section{PROJECTS AND RESEARCH}
%═══════════════════════════════════════
\begin{twocolentry}{Jul 2026}
\textbf{NYC Mobility Forecasting} [Python, PySpark, Delta Lake, scikit-learn, FastAPI]
\end{twocolentry}
\begin{highlights}
  \item Built a data acquisition and analysis system in \textbf{Python} ingesting trip, weather, and geospatial data through a Spark medallion pipeline into a \textbf{2.3M-row} hourly dataset across \textbf{263 zones}; two-stage Poisson hurdle model cut RMSE \textbf{43\%} and MAE \textbf{53\%} over baseline, served via a FastAPI REST API and dashboard.
\end{highlights}

\vspace{0.12 cm}

\begin{twocolentry}{Jul 2026}
\textbf{Spotify Analytics Platform} [Python, FastAPI, pandas, PostgreSQL, React]
\end{twocolentry}
\begin{highlights}
  \item Engineered a 4-stage \textbf{Python} ETL pipeline processing \textbf{71K+} streaming events with rate limiting, exponential-backoff retries, checkpointing, and validation; designed a PostgreSQL analytics schema with materialized views and \textbf{16 SQL RPC functions} for low-latency queries.
\end{highlights}

%═══════════════════════════════════════
\section{EXTRACURRICULAR ACTIVITIES}
%═══════════════════════════════════════
\begin{onecolentry}
\begin{highlights}
  \item \textbf{Finalist, RevolutionUC 2025 Hackathon} (300 participants, 31 schools) \,|\, \textbf{SAFe Practitioner 5.1} (Mar 2023).
\end{highlights}
\end{onecolentry}

\end{document}
